\documentclass[letterpaper]{article}

% -----------------------------------------------------------------------------
% Compile-safe AAAI wrapper
% -----------------------------------------------------------------------------
% If the official AAAI style file (aaai26.sty) is present in the project, this
% document will automatically use it. If not, it falls back to a standard article
% layout so the .tex file still compiles and produces a PDF. For final AAAI
% submission, upload aaai26.sty from the official AAAI author kit and compile
% again; the official style will then be used automatically.
\newif\ifuseaaai
\IfFileExists{aaai26.sty}{\useaaaitrue}{\useaaaifalse}

\ifuseaaai
  \usepackage[submission]{aaai26}
\else
  \usepackage[margin=0.9in]{geometry}
  \newcommand{\affiliations}[1]{\date{#1}}
\fi

\usepackage{times}
\usepackage{helvet}
\usepackage{courier}
\usepackage[hyphens]{url}
\usepackage{graphicx}
\urlstyle{rm}
\def\UrlFont{\rm}
\usepackage{natbib}
\usepackage{caption}
\usepackage{booktabs}
\usepackage{amsmath,amssymb,amsfonts}
\usepackage{amsthm}
\usepackage{multirow}
\usepackage{algorithm}
\usepackage{algorithmic}
\usepackage{xcolor}
\frenchspacing
\setlength{\pdfpagewidth}{8.5in}
\setlength{\pdfpageheight}{11in}

\newtheorem{theorem}{Theorem}
\newtheorem{proposition}{Proposition}
\newtheorem{remark}{Remark}

\title{InfoBridge: Entropy-Gated Relational Self-Distillation for Compact Text Embeddings}

\begin{document}

\maketitle

\begin{abstract}
Large text embedding models provide strong semantic representations but are costly to deploy. Distillation offers a practical way to compress them into compact encoders, yet layer-wise self-distillation for embedding models is still not fully understood. Existing relational self-distillation objectives often minimize the Frobenius distance between adjacent-layer relation matrices, thereby encouraging lower layers to copy the pairwise geometry of upper layers. We argue that this can be overly restrictive: lower layers should receive semantic guidance from upper layers without being forced to reproduce their entire relational structure. We propose \textsc{InfoBridge}, an entropy-gated relational self-distillation objective for compact text embeddings. The method converts layer-wise relation matrices into neighborhood distributions and uses entropy as a confidence signal to transfer only reliable semantic relations. We provide a focused theoretical analysis showing that hard relational matching uniformly penalizes all pairwise relations and can tie adjacent-layer geometries, while the proposed objective yields confidence-controlled gradients. The resulting proposal is intentionally simple: upper student layers are anchored to cached teacher embeddings, and lower layers are guided by selective self-distillation. This design avoids over-engineering while directly addressing the limitation of hard geometry matching in compact embedding distillation.
\end{abstract}

\section{Introduction}

Text embedding models are central to retrieval-augmented generation, semantic search, clustering, reranking, and many downstream NLP applications. Recent large embedding models achieve strong performance by combining large-scale pretraining, instruction tuning, and high-dimensional representation spaces. However, their deployment cost remains substantial, especially when embeddings must be computed at scale or served under latency constraints. Knowledge distillation (KD) is therefore a natural solution: a compact student encoder is trained to inherit the semantic behavior of a larger teacher model.

Embedding distillation differs from conventional classification distillation. In classification, the teacher usually provides a probability distribution over labels, and the student is optimized to match that distribution. In sentence embedding, the output is a continuous representation whose quality is determined by its geometry: semantic neighbors should be close, unrelated samples should be separated, and the representation space should transfer across domains. Consequently, a good distillation objective should preserve semantic structure rather than merely match individual vectors.

Recent work has explored efficient embedding distillation with only cached teacher sentence embeddings. A particularly relevant direction is to anchor the upper layers of a student to the teacher output while using self-distillation to propagate semantic structure to lower layers. TALAS is a recent example of this design: it uses the teacher's final embedding as an upper-layer anchor and introduces layer-aligned self-distillation by matching relation matrices between adjacent student layers. This avoids teacher hidden states and provides a practical way to bridge the capacity gap between large embedding teachers and compact students.

However, the self-distillation objective in this family of methods raises an important question: \emph{should a lower student layer copy the full relational geometry of the layer above it?} A typical relational self-distillation loss minimizes
\begin{equation}
    \sum_{l=1}^{L-1}\|R_{l+1}-R_l\|_F^2,
\end{equation}
where $R_l$ is a batch-wise relation matrix at layer $l$. This objective is intuitive and easy to implement, but it implicitly treats every pairwise relation as equally reliable and equally transferable. This may be too strong for compact embedding models. Lower layers often encode lexical and syntactic information, while upper layers encode more abstract semantic neighborhoods. Forcing lower layers to reproduce the full geometry of upper layers can reduce layer diversity and may transfer ambiguous or noisy relations.

This proposal studies a focused problem:
\begin{quote}
\emph{How can semantic information be transferred across student layers without forcing lower layers to clone the full upper-layer geometry?}
\end{quote}

We propose \textsc{InfoBridge}, an entropy-gated relational self-distillation objective. Instead of matching relation matrices directly, \textsc{InfoBridge} converts each row of a relation matrix into a neighborhood distribution. The entropy of the upper-layer distribution is used as a confidence signal: sharp neighborhoods provide stronger supervision, while high-entropy ambiguous neighborhoods are down-weighted. This reframes layer self-distillation from hard geometry matching to selective semantic relation transfer.

The intended contribution is deliberately narrow and defensible. We do not claim to solve knowledge distillation in general. We focus on adjacent-layer self-distillation for compact text embeddings and make three contributions:
\begin{enumerate}
    \item We identify a limitation of hard relational self-distillation: it uniformly penalizes all pairwise relations and can tie adjacent-layer geometries too strongly.
    \item We propose an entropy-gated relational objective that selectively transfers reliable semantic neighborhoods from upper layers to lower layers.
    \item We provide theoretical and empirical analysis plans to verify that selective transfer improves compact embedding distillation while preserving layer diversity.
\end{enumerate}

\section{Related Work}

\paragraph{Transformer Knowledge Distillation.}
Knowledge distillation was introduced as a way to transfer knowledge from a large teacher to a smaller student model \cite{hinton2015distilling}. In NLP, many methods distill intermediate Transformer representations rather than only final outputs. Patient Knowledge Distillation aligns student layers with selected teacher layers to compress BERT \cite{sun2019patient}. TinyBERT distills multiple components, including attention distributions and hidden states, in both general and task-specific stages \cite{jiao2020tinybert}. MiniLM focuses on self-attention relation distillation, while MobileBERT introduces a bottleneck architecture with carefully designed teacher-student transfer \cite{wang2020minilm,sun2020mobilebert}. These methods show that intermediate structure is important, but they often require teacher internal states. Such requirements are costly when the teacher is a large LLM-based embedding model or is only accessible through inference.

\paragraph{Relational Knowledge Distillation.}
Relational Knowledge Distillation (RKD) transfers relations among data samples instead of matching individual features coordinate-wise \cite{park2019relational}. This is especially relevant for embedding models, where neighborhood structure matters more than exact vector coordinates. Relational objectives are attractive because they are invariant to some coordinate-level differences and better reflect the geometry used in retrieval and semantic similarity. Nevertheless, common relational losses, such as Frobenius or MSE matching of relation matrices, still treat all pairwise relations uniformly. They do not distinguish reliable semantic neighborhoods from ambiguous or noisy relations.

\paragraph{Contrastive and Information-Theoretic Distillation.}
Contrastive Representation Distillation (CRD) argues that output-level KD may miss structural information in teacher representations and uses a contrastive objective to transfer richer information \cite{tian2019crd}. Variational Information Distillation (VID) interprets KD as maximizing mutual information between teacher and student representations \cite{ahn2019vid}. These approaches motivate an information-theoretic view of distillation, but most of them focus on teacher-student representation transfer rather than adjacent-layer self-distillation inside a compact embedding model. Our proposal uses this perspective more narrowly: the goal is not to estimate mutual information as a universal KD objective, but to design a practical confidence-aware self-distillation loss for layer-wise semantic transfer.

\paragraph{Sentence Embedding Learning and Distillation.}
Sentence embedding learning has been strongly influenced by contrastive objectives such as SimCSE \cite{gao2021simcse}, which improve alignment and uniformity of embedding spaces. MTEB further highlights the need to evaluate embeddings across diverse tasks, including retrieval, reranking, classification, clustering, and semantic textual similarity \cite{muennighoff2023mteb}. Recent sentence embedding distillation methods such as DistillCSE, EMO, Jasper and Stella, and ESE explore contrastive, relational, optimal transport, or compression-based mechanisms for transferring semantic representations \cite{xu2023distillcse,truong2025emo,zhang2025jasper,li2025ese}. TALAS is particularly relevant because it combines cached teacher embedding anchoring with adjacent-layer relational self-distillation. We use TALAS as a recent prior work demonstrating the usefulness of layer self-distillation, but our work addresses a more general question: what kind of self-distillation objective is appropriate when lower layers should receive semantic guidance without copying the full upper-layer geometry?

\section{Problem Formulation}

Let $T$ be a teacher embedding model and $S$ be a compact student encoder with $L$ layers. Given a mini-batch of $N$ sentences
\begin{equation}
    B=\{x_i\}_{i=1}^{N},
\end{equation}
the teacher produces final sentence embeddings
\begin{equation}
    E^T = T(B) \in \mathbb{R}^{N\times d_T}.
\end{equation}
The student produces layer-wise sentence embeddings
\begin{equation}
    E_l^S = S_l(B) \in \mathbb{R}^{N\times d_l},\quad l=1,\ldots,L.
\end{equation}
The teacher embeddings are pre-computed and cached. During student training, no teacher hidden states or online teacher forward passes are required.

The overall training objective is kept simple:
\begin{equation}
    \mathcal{L}=\mathcal{L}_{\mathrm{anchor}}+\lambda\mathcal{L}_{\mathrm{IBridge}}.
\end{equation}
The first term anchors upper student layers to the teacher output. The second term is the proposed selective relational self-distillation loss.

\section{Method: InfoBridge}

\subsection{Upper-Layer Teacher Anchoring}

We use cached teacher embeddings to supervise only the upper $K$ student layers:
\begin{equation}
    \mathcal{L}_{\mathrm{anchor}}
    =
    \frac{1}{K}
    \sum_{l=L-K+1}^{L}
    D_{\cos}(E_l^S W_l,E^T),
\end{equation}
where $W_l\in\mathbb{R}^{d_l\times d_T}$ is a learnable projection matrix and $D_{\cos}$ is cosine distance averaged over the batch. This component follows the common intuition that upper layers are better suited to approximate teacher-level semantic abstraction. It is not the main novelty; it provides the semantic source that self-distillation propagates downward.

\subsection{Layer-Wise Relation Matrices}

For each layer, we first normalize the embeddings row-wise:
\begin{equation}
    \bar{E}_l^S = \mathrm{Normalize}(E_l^S).
\end{equation}
We then compute the batch-wise relation matrix:
\begin{equation}
    R_l = \bar{E}_l^S(\bar{E}_l^S)^\top \in \mathbb{R}^{N\times N}.
\end{equation}
The entry $(R_l)_{ij}$ is the cosine similarity between samples $i$ and $j$ at layer $l$. Hard relational self-distillation directly matches $R_l$ and $R_{l+1}$. In contrast, \textsc{InfoBridge} treats each row of $R_l$ as a local semantic neighborhood.

\subsection{Neighborhood Distributions}

For each sample $i$, we remove the diagonal self-similarity and convert the remaining similarities into a distribution:
\begin{equation}
    p_i^{l+1}
    =
    \mathrm{softmax}\left(\frac{R_{i,\neg i}^{l+1}}{\tau_t}\right),
\end{equation}
\begin{equation}
    q_i^l
    =
    \mathrm{softmax}\left(\frac{R_{i,\neg i}^{l}}{\tau_s}\right),
\end{equation}
where $R_{i,\neg i}^{l}$ denotes row $i$ with the diagonal entry removed, and $\tau_t,\tau_s>0$ are temperatures. The distribution $p_i^{l+1}$ describes the upper-layer neighborhood of sample $i$, while $q_i^l$ describes the lower-layer neighborhood.

This distributional view is less restrictive than full matrix matching. It focuses on relative neighborhood structure rather than exact similarity magnitudes for all pairs.

\subsection{Entropy-Gated Self-Distillation}

A key observation is that not every upper-layer neighborhood is equally informative. If $p_i^{l+1}$ is sharp, the upper layer has identified clear semantic neighbors for $x_i$. If it is close to uniform, the semantic neighborhood is ambiguous. We therefore define an entropy-based confidence score:
\begin{equation}
    \alpha_i^{l+1}
    =
    1-\frac{H(p_i^{l+1})}{\log(N-1)},
\end{equation}
where
\begin{equation}
    H(p_i^{l+1})=-\sum_j p_{ij}^{l+1}\log p_{ij}^{l+1}.
\end{equation}
Since $0\leq H(p_i^{l+1})\leq \log(N-1)$, we have $0\leq \alpha_i^{l+1}\leq 1$.

The proposed loss for adjacent layers $l$ and $l+1$ is
\begin{equation}
    \mathcal{L}_{\mathrm{IBridge}}^{(l)}
    =
    \frac{1}{N}
    \sum_{i=1}^{N}
    \alpha_i^{l+1}
    \mathrm{KL}\left(
    \mathrm{sg}(p_i^{l+1})\|q_i^l
    \right),
\end{equation}
where $\mathrm{sg}(\cdot)$ denotes stop-gradient. The full self-distillation loss is
\begin{equation}
    \mathcal{L}_{\mathrm{IBridge}}
    =
    \frac{1}{L-1}
    \sum_{l=1}^{L-1}
    \mathcal{L}_{\mathrm{IBridge}}^{(l)}.
\end{equation}

The role of the entropy gate is simple: confident upper-layer neighborhoods are transferred strongly; uncertain neighborhoods are transferred weakly. This avoids treating all batch relations as equally reliable.

\begin{algorithm}[t]
\caption{Training with InfoBridge}
\begin{algorithmic}[1]
\REQUIRE Training corpus $\mathcal{D}$, cached teacher embeddings $E^T$, student $S$, number of anchored layers $K$, weight $\lambda$.
\FOR{each mini-batch $B\subset\mathcal{D}$}
    \STATE Retrieve cached teacher embeddings $E^T_B$.
    \STATE Compute student layer embeddings $\{E_l^S\}_{l=1}^L$.
    \STATE Compute $\mathcal{L}_{\mathrm{anchor}}$ on the upper $K$ layers.
    \FOR{$l=1$ to $L-1$}
        \STATE Compute relation matrices $R_l$ and $R_{l+1}$.
        \STATE Convert each row to distributions $q_i^l$ and $p_i^{l+1}$.
        \STATE Compute entropy gates $\alpha_i^{l+1}$.
        \STATE Accumulate $\mathcal{L}_{\mathrm{IBridge}}^{(l)}$.
    \ENDFOR
    \STATE Update $S$ using $\mathcal{L}_{\mathrm{anchor}}+\lambda\mathcal{L}_{\mathrm{IBridge}}$.
\ENDFOR
\end{algorithmic}
\end{algorithm}

\section{Theoretical Analysis}

This section provides focused theoretical support for the proposed objective. The goal is not to prove a broad generalization guarantee. Instead, we formalize why hard relational matching can be restrictive and why entropy-gated transfer is a principled relaxation.

\subsection{Hard Relational Matching Uniformly Penalizes All Relations}

\begin{proposition}[Uniform pairwise penalization]
Let $R_A,R_B\in\mathbb{R}^{N\times N}$ be two relation matrices and consider
\begin{equation}
    \mathcal{L}_{\mathrm{rel}}=\|R_A-R_B\|_F^2.
\end{equation}
Then for every pair $(i,j)$,
\begin{equation}
    \frac{\partial \mathcal{L}_{\mathrm{rel}}}{\partial (R_A)_{ij}}
    =
    2((R_A)_{ij}-(R_B)_{ij}).
\end{equation}
\end{proposition}

\begin{proof}
By definition,
\begin{equation}
    \mathcal{L}_{\mathrm{rel}}
    =
    \sum_{i,j}((R_A)_{ij}-(R_B)_{ij})^2.
\end{equation}
Differentiating with respect to $(R_A)_{ij}$ gives the result.
\end{proof}

\paragraph{Implication.}
The gradient depends only on the residual for each pair. There is no term reflecting whether a relation is semantically reliable, ambiguous, or noisy. Thus, Frobenius relational matching performs uniform pairwise transfer. This is acceptable when all relations are trusted, but it can be restrictive for layer self-distillation where different samples may have different neighborhood certainty.

\subsection{Hard Adjacent-Layer Matching Can Tie Layer Geometries}

\begin{theorem}[Geometry tying under adjacent-layer relational matching]
Let $R_1,\ldots,R_L$ be relation matrices for $L$ student layers. Suppose the average adjacent-layer matching loss satisfies
\begin{equation}
    \frac{1}{L-1}\sum_{l=1}^{L-1}\|R_{l+1}-R_l\|_F^2 \leq \epsilon.
\end{equation}
Then for any $1\leq a<b\leq L$,
\begin{equation}
    \|R_b-R_a\|_F
    \leq
    (b-a)\sqrt{(L-1)\epsilon}.
\end{equation}
In particular, if $\epsilon=0$, then $R_1=R_2=\cdots=R_L$.
\end{theorem}

\begin{proof}
Since
\begin{equation}
    \sum_{l=1}^{L-1}\|R_{l+1}-R_l\|_F^2 \leq (L-1)\epsilon,
\end{equation}
we have for every $l$,
\begin{equation}
    \|R_{l+1}-R_l\|_F \leq \sqrt{(L-1)\epsilon}.
\end{equation}
Using the triangle inequality,
\begin{align}
    \|R_b-R_a\|_F
    &=
    \left\|\sum_{l=a}^{b-1}(R_{l+1}-R_l)\right\|_F \\
    &\leq
    \sum_{l=a}^{b-1}\|R_{l+1}-R_l\|_F \\
    &\leq
    (b-a)\sqrt{(L-1)\epsilon}.
\end{align}
If $\epsilon=0$, every adjacent difference is zero, hence all relation matrices are identical.
\end{proof}

\paragraph{Implication.}
When optimized strongly, hard adjacent-layer relational matching encourages the relational geometries of multiple layers to become close. This formally supports the concern that such objectives may reduce layer-wise geometric diversity. The theorem does not claim that this always harms performance; rather, it identifies a structural bias of the loss.

\subsection{Entropy-Gated KL Gives Confidence-Controlled Gradients}

\begin{theorem}[Confidence-controlled gradient]
Let $p\in\Delta^{N-2}$ be an upper-layer neighborhood distribution, $q=\mathrm{softmax}(s)$ be a lower-layer neighborhood distribution, and
\begin{equation}
    \mathcal{L}(s)=\alpha\,\mathrm{KL}(p\|q),
\end{equation}
where $\alpha\in[0,1]$ is independent of $s$. Then
\begin{equation}
    \nabla_s\mathcal{L}=\alpha(q-p),
\end{equation}
and
\begin{equation}
    \|\nabla_s\mathcal{L}\|_2\leq \alpha\sqrt{2}.
\end{equation}
For the proposed entropy gate $\alpha=1-H(p)/\log(N-1)$, if $p$ is uniform, then $\alpha=0$ and $\nabla_s\mathcal{L}=0$.
\end{theorem}

\begin{proof}
For $q=\mathrm{softmax}(s)$,
\begin{equation}
    \nabla_s\mathrm{KL}(p\|q)=q-p.
\end{equation}
Multiplying by $\alpha$ gives the first result. Since $p$ and $q$ are probability vectors, the maximum Euclidean distance between them is at most $\sqrt{2}$, achieved by two different one-hot distributions. Hence $\|\alpha(q-p)\|_2\leq\alpha\sqrt{2}$. If $p$ is uniform, $H(p)=\log(N-1)$ and therefore $\alpha=0$.
\end{proof}

\paragraph{Implication.}
The proposed loss automatically reduces supervision from high-entropy neighborhoods. Unlike Frobenius matching, it does not force the student to imitate uncertain relations with the same strength as confident ones.

\subsection{Relation to Information Transfer}

The proposed objective can be interpreted as selective information transfer rather than full geometry cloning. Each row-wise distribution $p_i^{l+1}$ summarizes the semantic neighborhood information available at layer $l+1$. The KL term transfers this neighborhood information to layer $l$, while the entropy gate controls the amount of transfer based on uncertainty. This is a practical information-theoretic perspective, not a claim that the method exactly optimizes mutual information. The formulation intentionally avoids additional latent-variable bottlenecks or complex estimators so that the method remains simple and easy to evaluate.

\section{Experimental Design}

\subsection{Main Question}

The main experimental question is:
\begin{quote}
Does entropy-gated relational self-distillation improve compact text embeddings while preserving layer diversity better than hard self-distillation losses?
\end{quote}
To answer this, the primary experiments should isolate the self-distillation objective. All variants should use the same teacher, student, data, teacher-anchor loss, optimizer, batch size, and training steps. Only the self-distillation loss should change.

\subsection{Baselines}

The clean comparison uses the same objective form
\begin{equation}
    \mathcal{L}=\mathcal{L}_{\mathrm{anchor}}+\lambda\mathcal{L}_{\mathrm{SD}},
\end{equation}
with different choices of $\mathcal{L}_{\mathrm{SD}}$:
\begin{itemize}
    \item \textbf{No SD}: teacher anchor only.
    \item \textbf{Hidden MSE SD}: point-wise adjacent-layer representation matching.
    \item \textbf{Cosine SD}: normalized point-wise adjacent-layer matching.
    \item \textbf{Frobenius Relation SD}: hard matching of adjacent-layer relation matrices.
    \item \textbf{KL Relation SD}: row-wise neighborhood distribution matching without entropy gates.
    \item \textbf{InfoNCE SD}: contrastive adjacent-layer transfer.
    \item \textbf{InfoBridge}: entropy-gated relational self-distillation.
\end{itemize}
External baselines may include representative sentence embedding distillation methods, but the most important comparison is the controlled loss-level comparison above.

\subsection{Teacher--Student Settings}

We propose three settings covering different capacity gaps:
\begin{itemize}
    \item Qwen3-Embedding-0.6B $\rightarrow$ MiniLMv2-H384.
    \item BGE-M3 $\rightarrow$ MiniLMv2-H768.
    \item Qwen3-Embedding-4B $\rightarrow$ BERT-base.
\end{itemize}
These settings are appropriate for testing whether selective self-distillation is especially useful when the semantic teacher is much stronger than the compact student.

\subsection{Benchmarks}

The minimum evaluation should include semantic textual similarity, pair classification, and classification tasks. A stronger submission should include a representative MTEB subset, especially retrieval and reranking tasks, because those settings directly depend on embedding neighborhood quality. Results should be reported both in-domain and out-of-domain when possible.

\subsection{Diagnostics}

Performance alone is not enough to support the proposed hypothesis. We therefore plan the following analyses:
\begin{itemize}
    \item \textbf{Layer diversity}: measure CKA or representational similarity between student layers.
    \item \textbf{Relation distance}: track $\|R_{l+1}-R_l\|_F$ across layers.
    \item \textbf{Top-$k$ neighborhood preservation}: measure overlap between nearest-neighbor sets across layers.
    \item \textbf{Gradient norm by layer}: measure $\|\nabla_{H_l}\mathcal{L}_{\mathrm{SD}}\|$ to test whether hard matching creates stronger pressure on shallow layers.
    \item \textbf{Entropy-confidence visualization}: inspect the distribution of $\alpha_i$ and analyze examples with low and high entropy.
    \item \textbf{Embedding geometry}: evaluate anisotropy, alignment, and uniformity to check whether the final embedding space remains well-structured.
\end{itemize}
These diagnostics directly test whether the proposed objective transfers semantic structure selectively rather than merely improving benchmark scores.

\subsection{Ablations}

The most important ablations are:
\begin{itemize}
    \item without entropy gating, i.e., KL relation matching only;
    \item different entropy-to-confidence functions;
    \item different temperatures $\tau_t$ and $\tau_s$;
    \item different self-distillation weights $\lambda$;
    \item different numbers of anchored upper layers $K$;
    \item top-$k$ neighborhood distribution versus full row distribution;
    \item with and without additional contrastive regularization.
\end{itemize}
The expected decisive comparison is Frobenius Relation SD versus KL Relation SD versus InfoBridge. This tests whether distributional matching and entropy-based selectivity are both necessary.

\section{Expected Results and Interpretation}

We expect InfoBridge to improve over hard relational self-distillation particularly on out-of-domain or retrieval-oriented evaluations, where preserving transferable semantic neighborhoods is important. We also expect InfoBridge to maintain lower layer similarity than hard matching in CKA analysis, reduce unnecessary gradient pressure on shallow layers, and preserve top-$k$ neighborhoods without forcing the entire relation matrix to match.

These expectations are intentionally modest. The claim is not that entropy-gated self-distillation will dominate every KD method or every benchmark. The claim is that, for adjacent-layer self-distillation in compact text embeddings, selective relation transfer is a better inductive bias than uniform full-geometry matching.

\section{Limitations}

This proposal has several limitations. First, entropy is a simple confidence proxy and may be affected by batch composition. Second, row-wise neighborhood distributions depend on batch size, so small batches may provide weaker relational signals. Third, the method focuses on text embedding encoders and may not directly transfer to generative language model distillation. Fourth, the theory characterizes structural properties of the losses rather than providing a full generalization guarantee. These limitations are acceptable for the intended scope: a focused study of selective adjacent-layer self-distillation for compact embedding models.

\section{Conclusion}

This proposal reframes layer self-distillation for compact text embeddings from hard geometry matching to selective semantic relation transfer. The method is intentionally simple: use cached teacher embeddings to anchor upper student layers, then use entropy-gated relational self-distillation to propagate only confident semantic neighborhoods downward. Theoretical analysis identifies why hard relational matching can be overly restrictive, and the proposed experiments are designed to verify whether selective transfer improves performance while preserving layer diversity. The resulting work is positioned as a focused contribution to embedding distillation, with TALAS and relational KD serving as important prior work rather than as the sole target of improvement.

\section*{Acknowledgments}
This section is omitted for anonymous submission.

\begin{thebibliography}{99}

\bibitem[Ahn et al.(2019)]{ahn2019vid}
Ahn, S.; Hu, S. X.; Damianou, A.; Lawrence, N. D.; and Dai, Z. 2019.
Variational Information Distillation for Knowledge Transfer.
In \emph{Proceedings of CVPR}.

\bibitem[Gao, Yao, and Chen(2021)]{gao2021simcse}
Gao, T.; Yao, X.; and Chen, D. 2021.
SimCSE: Simple Contrastive Learning of Sentence Embeddings.
In \emph{Proceedings of EMNLP}.

\bibitem[Hinton, Vinyals, and Dean(2015)]{hinton2015distilling}
Hinton, G.; Vinyals, O.; and Dean, J. 2015.
Distilling the Knowledge in a Neural Network.
\emph{arXiv preprint arXiv:1503.02531}.

\bibitem[Jiao et al.(2020)]{jiao2020tinybert}
Jiao, X.; Yin, Y.; Shang, L.; Jiang, X.; Chen, X.; Li, L.; Wang, F.; and Liu, Q. 2020.
TinyBERT: Distilling BERT for Natural Language Understanding.
In \emph{Findings of EMNLP}.

\bibitem[Li et al.(2025)]{li2025ese}
Li, X.; Li, Z.; Li, J.; Xie, H.; and Li, Q. 2025.
ESE: Espresso Sentence Embeddings.
In \emph{Proceedings of ICLR}.

\bibitem[Muennighoff et al.(2023)]{muennighoff2023mteb}
Muennighoff, N.; Tazi, N.; Magne, L.; and Reimers, N. 2023.
MTEB: Massive Text Embedding Benchmark.
In \emph{Proceedings of EACL}.

\bibitem[Park et al.(2019)]{park2019relational}
Park, W.; Kim, D.; Lu, Y.; and Cho, M. 2019.
Relational Knowledge Distillation.
In \emph{Proceedings of CVPR}.

\bibitem[Sun et al.(2019)]{sun2019patient}
Sun, S.; Cheng, Y.; Gan, Z.; and Liu, J. 2019.
Patient Knowledge Distillation for BERT Model Compression.
In \emph{Proceedings of EMNLP-IJCNLP}.

\bibitem[Sun et al.(2020)]{sun2020mobilebert}
Sun, Z.; Yu, H.; Song, X.; Liu, R.; Yang, Y.; and Zhou, D. 2020.
MobileBERT: a Compact Task-Agnostic BERT for Resource-Limited Devices.
In \emph{Proceedings of ACL}.

\bibitem[Tian, Krishnan, and Isola(2020)]{tian2019crd}
Tian, Y.; Krishnan, D.; and Isola, P. 2020.
Contrastive Representation Distillation.
In \emph{Proceedings of ICLR}.

\bibitem[Truong et al.(2025)]{truong2025emo}
Truong, M.-P.; Vu, H. A.; Vu, T.; Diep, N. T. N.; Van, L. N.; Nguyen, T. H.; and Le, T. 2025.
EMO: Embedding Model Distillation via Intra-Model Relation and Optimal Transport Alignments.
In \emph{Proceedings of EMNLP}.

\bibitem[Wang et al.(2020)]{wang2020minilm}
Wang, W.; Wei, F.; Dong, L.; Bao, H.; Yang, N.; and Zhou, M. 2020.
MiniLM: Deep Self-Attention Distillation for Task-Agnostic Compression of Pre-Trained Transformers.
In \emph{Proceedings of NeurIPS}.

\bibitem[Xu et al.(2023)]{xu2023distillcse}
Xu, J.; Shao, W.; Chen, L.; and Liu, L. 2023.
Distilled Contrastive Learning for Sentence Embeddings.
In \emph{Findings of EMNLP}.

\bibitem[Zhang et al.(2025)]{zhang2025jasper}
Zhang, D.; Li, J.; Zeng, Z.; and Wang, F. 2025.
Jasper and Stella: Distillation of SOTA Embedding Models.
\emph{arXiv preprint}.

\bibitem[Anonymous(2026)]{anonymous2026talas}
Anonymous. 2026.
TALAS: Teacher-Anchored Layer Alignment with Adaptive Sharpness-Aware Minimization for Embedding Distillation.
Anonymous ACL submission.

\end{thebibliography}

\end{document}
